% --- Archon physics formalization source begin ---
% archon:physics
% archon:covers IPhO2026Problems/problem_IPhO_2026_4_A_5.lean
% archon:source-report reports/ipho_2026/problem_IPhO_2026_4_A_5.source.json
% archon:problem-id IPhO_2026_4
% archon:part-id A.5
% archon:previous-part IPhO_2026_4_A_3
% archon:previous-part-policy natural_language_prerequisite_only

\chapter{Physics problem IPhO\_2026\_4\_A\_5}
\label{ch:IPhO2026Problems_problem_IPhO_2026_4_A_5}

\paragraph{Problem source.}
The experimental apparatus contains a sealed air column (CA) in the inner
cylinder.  Propylene glycol is introduced to h = 4.5 cm so the air volume is
fixed.  Use the cylinder dimensions in Figure 17, ambient air density
rho\_a = 1.12 kg/m\textasciicircum{}3, and the ideal-gas law P*V = n*R*T.  The outer-cylinder
water bath is heated while pressure and temperature are recorded.

Current subquestion:
Determine the constant-volume thermal pressure coefficient beta\_0 = (1/P\_0)*(Delta P/Delta T).

\paragraph{Current subquestion.}
Determine the constant-volume thermal pressure coefficient beta\_0 = (1/P\_0)*(Delta P/Delta T).

\paragraph{Recorded answer/context.}
Official sample: beta\_0 = 0.0034 +/- 0.0007 K\textasciicircum{}(-1); ideal-gas reference 1/273.15 K = 0.0037 K\textasciicircum{}(-1).

\paragraph{Figure/image path.}
/root/proposal\_for\_physic/science-mango/ipho\_2026\_source/image/E1\_page-9.png

\paragraph{Reusable previous-part conclusions.}
\begin{itemize}
\item Source A.3. Question: Plot pressure as a function of temperature from A2. Reusable conclusions: The expected isochoric ideal-gas plot is linear: P is proportional to absolute T. Policy: natural\_language\_prerequisite\_only; do\_not\_import\_Lean\_output
\end{itemize}

\paragraph{Formalization target.}
create a compiling Lean file with sorry bodies at `IPhO2026Problems/problem\_IPhO\_2026\_4\_A\_5.lean`. The Lean declarations must preserve the physical quantities, dimensions or dimensional roles, figure labels, governing-law hypotheses, and final relation expressed by this problem.
Use Mathlib/Physlib names found through LeanExplore where available. If a physics API is missing, introduce faithful local abstractions rather than scalar placeholder aliases.

\paragraph{Iteration 002 redraft contract.}
Keep the experimental coefficient separate from the ideal-gas comparison.
Introduce a measured reference pressure and a pressure-versus-temperature
secant or regression slope, each with an absolute uncertainty covering the
underlying quantity.  Prove the general quotient formula
\(\beta_0=s/P_0\) and propagate the two input intervals through this positive
quotient.  The official solution data include the fitted central slope
\(0.3229\ {\rm kPa\,K^{-1}}\) and reference pressure
\(P_0=101.3\) kPa; use a regression/error certificate, not membership in a
fixed output band, as the experimental premise.  Separately derive the ideal
reference \(1/273.15\ {\rm K^{-1}}\) and its \(0.0037\) rounding.

\begin{theorem}[Physics formalization target]
  \label{thm:physics:IPhO_2026_4_A_5:target}
  \lean{IPhO2026Problems.Problem4A5.IPhO2026_4_A_5_thermalPressureCoefficient}
  \uses{def:physics:IPhO_2026_4_A_5:aux001, def:physics:IPhO_2026_4_A_5:aux002, def:physics:IPhO_2026_4_A_5:aux003, def:physics:IPhO_2026_4_A_5:aux004, def:physics:IPhO_2026_4_A_5:aux005, def:physics:IPhO_2026_4_A_5:aux006, def:physics:IPhO_2026_4_A_5:aux007, def:physics:IPhO_2026_4_A_5:aux008, def:physics:IPhO_2026_4_A_5:aux009, def:physics:IPhO_2026_4_A_5:aux010, def:physics:IPhO_2026_4_A_5:aux011, def:physics:IPhO_2026_4_A_5:aux012, def:physics:IPhO_2026_4_A_5:aux013, def:physics:IPhO_2026_4_A_5:aux014, def:physics:IPhO_2026_4_A_5:aux015, lem:physics:IPhO_2026_4_A_5:aux016, lem:physics:IPhO_2026_4_A_5:aux017}
  Determine the experimental constant-volume pressure coefficient from the
  measured positive regression slope and reference pressure, propagate their
  uncertainties to a coefficient interval, and keep this result distinct
  from the ideal-gas value \(1/273.15\ {\rm K^{-1}}\).
\end{theorem}
\begin{proof}
  By definition the pressure change along the fitted line is
  \(s\,\Delta T\), so \(\beta_0=(1/P_0)(\Delta P/\Delta T)=s/P_0\).
  Positivity makes this quotient increasing in \(s\) and decreasing in
  \(P_0\), giving the endpoint error interval.  Instantiate the regression
  certificate for the experimental report.  For the ideal comparison,
  isochoric \(PV=nRT\) gives \(s=P_0/T_0\), hence
  \(\beta_0=1/T_0\).
\end{proof}
\section*{Formal declaration index}
The following blocks record the typed carriers and bridge declarations used by the target.  Their dependency lists follow the declaration-level model in the covered Lean file.

\begin{definition}[Declaration ApparatusLabel]
  \label{def:physics:IPhO_2026_4_A_5:aux001}
  \lean{IPhO2026Problems.Problem4A5.ApparatusLabel}
  Labels used for the apparatus in Figure 17 and on the official source page.
\end{definition}
\begin{proof}
  This is the typed definition of the stated physical carrier; its fields or defining equation provide the claimed data.
\end{proof}

\begin{definition}[Declaration Figure17Geometry]
  \label{def:physics:IPhO_2026_4_A_5:aux002}
  \lean{IPhO2026Problems.Problem4A5.Figure17Geometry}
  The cylinder dimensions read from Figure 17, expressed as SI scalar readouts.
\end{definition}
\begin{proof}
  This is the typed definition of the stated physical carrier; its fields or defining equation provide the claimed data.
\end{proof}

\begin{definition}[Declaration IsochoricApparatus]
  \label{def:physics:IPhO_2026_4_A_5:aux003}
  \lean{IPhO2026Problems.Problem4A5.IsochoricApparatus}
  \uses{def:physics:IPhO_2026_4_A_5:aux002}
  The part of the experimental apparatus relevant to the sealed air column. Scalar fields are explicitly SI readouts; pressure itself is represented below by Physlib's dimensionful DimPressure.
\end{definition}
\begin{proof}
  This is the typed definition of the stated physical carrier; its fields or defining equation provide the claimed data.
\end{proof}

\begin{definition}[Declaration IsPreparedIsochoricApparatus]
  \label{def:physics:IPhO_2026_4_A_5:aux004}
  \lean{IPhO2026Problems.Problem4A5.IsPreparedIsochoricApparatus}
  \uses{def:physics:IPhO_2026_4_A_5:aux003}
  The Figure 17 geometry and the instructions h = 4.5 cm, ρₐ = 1.12 kg/m³, with both valves closed. The final equality is the cylindrical volume of CA remaining above the propylene glycol.
\end{definition}
\begin{proof}
  This is the typed definition of the stated physical carrier; its fields or defining equation provide the claimed data.
\end{proof}

\begin{definition}[Declaration pressureInPascals]
  \label{def:physics:IPhO_2026_4_A_5:aux005}
  \lean{IPhO2026Problems.Problem4A5.pressureInPascals}
  The numerical value, in pascals, of a dimensionful pressure in SI units.
\end{definition}
\begin{proof}
  This is the typed definition of the stated physical carrier; its fields or defining equation provide the claimed data.
\end{proof}

\begin{definition}[Declaration temperatureInKelvin]
  \label{def:physics:IPhO_2026_4_A_5:aux006}
  \lean{IPhO2026Problems.Problem4A5.temperatureInKelvin}
  The real readout of an absolute temperature. Throughout this experiment the arbitrary temperature unit carried by Temperature is fixed to kelvin.
\end{definition}
\begin{proof}
  This is the typed definition of the stated physical carrier; its fields or defining equation provide the claimed data.
\end{proof}

\begin{definition}[Declaration IsochoricHeatingRun]
  \label{def:physics:IPhO_2026_4_A_5:aux007}
  \lean{IPhO2026Problems.Problem4A5.IsochoricHeatingRun}
  \uses{def:physics:IPhO_2026_4_A_5:aux003}
  One constant-volume heating run. pressureAt is a physical, dimensionful pressure, while the amount of substance and gas constant are SI scalar readouts in mol and J mol⁻¹ K⁻¹.
\end{definition}
\begin{proof}
  This is the typed definition of the stated physical carrier; its fields or defining equation provide the claimed data.
\end{proof}

\begin{definition}[Declaration UsesStandardReferenceState]
  \label{def:physics:IPhO_2026_4_A_5:aux008}
  \lean{IPhO2026Problems.Problem4A5.UsesStandardReferenceState}
  \uses{def:physics:IPhO_2026_4_A_5:aux005, def:physics:IPhO_2026_4_A_5:aux006, def:physics:IPhO_2026_4_A_5:aux007}
  The system reference temperature is 273.15 K and its pressure is positive.
\end{definition}
\begin{proof}
  This is the typed definition of the stated physical carrier; its fields or defining equation provide the claimed data.
\end{proof}

\begin{definition}[Declaration IsHeatingBranch]
  \label{def:physics:IPhO_2026_4_A_5:aux009}
  \lean{IPhO2026Problems.Problem4A5.IsHeatingBranch}
  \uses{def:physics:IPhO_2026_4_A_5:aux005, def:physics:IPhO_2026_4_A_5:aux006, def:physics:IPhO_2026_4_A_5:aux007}
  Orientation information for the recorded branch: the second state is reached by heating, and the measured pressure increases along that branch.
\end{definition}
\begin{proof}
  This is the typed definition of the stated physical carrier; its fields or defining equation provide the claimed data.
\end{proof}

\begin{definition}[Declaration ObeysIsochoricIdealGasLaw]
  \label{def:physics:IPhO_2026_4_A_5:aux010}
  \lean{IPhO2026Problems.Problem4A5.ObeysIsochoricIdealGasLaw}
  \uses{def:physics:IPhO_2026_4_A_5:aux005, def:physics:IPhO_2026_4_A_5:aux006, def:physics:IPhO_2026_4_A_5:aux007}
  The governing ideal-gas law P V = n R T for the sealed air column. The one apparatus volume is used at every temperature, encoding the isochoric process.
\end{definition}
\begin{proof}
  This is the typed definition of the stated physical carrier; its fields or defining equation provide the claimed data.
\end{proof}

\begin{definition}[Declaration HasIsochoricPressureLinearity]
  \label{def:physics:IPhO_2026_4_A_5:aux011}
  \lean{IPhO2026Problems.Problem4A5.HasIsochoricPressureLinearity}
  \uses{def:physics:IPhO_2026_4_A_5:aux005, def:physics:IPhO_2026_4_A_5:aux006, def:physics:IPhO_2026_4_A_5:aux007, def:physics:IPhO_2026_4_A_5:aux009}
  Reusable content of part A.3: the pressure plot is proportional to absolute temperature. The quantified equation makes this a constraining interface.
\end{definition}
\begin{proof}
  This is the typed definition of the stated physical carrier; its fields or defining equation provide the claimed data.
\end{proof}

\begin{definition}[Declaration thermalPressureCoefficientPerKelvin]
  \label{def:physics:IPhO_2026_4_A_5:aux012}
  \lean{IPhO2026Problems.Problem4A5.thermalPressureCoefficientPerKelvin}
  \uses{def:physics:IPhO_2026_4_A_5:aux005, def:physics:IPhO_2026_4_A_5:aux006, def:physics:IPhO_2026_4_A_5:aux007, def:physics:IPhO_2026_4_A_5:aux009}
  The source definition β₀ = (1 / P₀) (ΔP / ΔT), with the reference-to-heated orientation fixed by IsHeatingBranch. Its dimensional role is inverse kelvin.
\end{definition}
\begin{proof}
  This is the typed definition of the stated physical carrier; its fields or defining equation provide the claimed data.
\end{proof}

\begin{definition}[Declaration Estimate]
  \label{def:physics:IPhO_2026_4_A_5:aux013}
  \lean{IPhO2026Problems.Problem4A5.Estimate}
  A scalar estimate together with its symmetric reported uncertainty.
\end{definition}
\begin{proof}
  This is the typed definition of the stated physical carrier; its fields or defining equation provide the claimed data.
\end{proof}

\begin{definition}[Declaration Estimate.Contains]
  \label{def:physics:IPhO_2026_4_A_5:aux014}
  \lean{IPhO2026Problems.Problem4A5.Estimate.Contains}
  \uses{def:physics:IPhO_2026_4_A_5:aux013}
  Membership in the closed symmetric uncertainty band of an estimate.
\end{definition}
\begin{proof}
  This is the typed definition of the stated physical carrier; its fields or defining equation provide the claimed data.
\end{proof}

\begin{definition}[Declaration officialCoefficientEstimatePerKelvin]
  \label{def:physics:IPhO_2026_4_A_5:aux015}
  \lean{IPhO2026Problems.Problem4A5.officialCoefficientEstimatePerKelvin}
  \uses{def:physics:IPhO_2026_4_A_5:aux013}
  The official experimental answer 0.0034 ± 0.0007 K⁻¹.
\end{definition}
\begin{proof}
  This is the typed definition of the stated physical carrier; its fields or defining equation provide the claimed data.
\end{proof}

\begin{lemma}[Declaration idealGasLaw\_implies\_isochoricPressureLinearity]
  \label{lem:physics:IPhO_2026_4_A_5:aux016}
  \lean{IPhO2026Problems.Problem4A5.idealGasLaw_implies_isochoricPressureLinearity}
  \uses{def:physics:IPhO_2026_4_A_5:aux004, def:physics:IPhO_2026_4_A_5:aux007, def:physics:IPhO_2026_4_A_5:aux010, def:physics:IPhO_2026_4_A_5:aux011}
  The ideal-gas law at fixed positive volume supplies the proportionality used in the pressure-versus-temperature plot of part A.3.
\end{lemma}
\begin{proof}
  Substitute the cited governing relations in order and use the stated positivity, orientation, and branch conditions to obtain the conclusion.
\end{proof}

\begin{lemma}[Declaration thermalPressureCoefficient\_eq\_inverse\_referenceTemperature]
  \label{lem:physics:IPhO_2026_4_A_5:aux017}
  \lean{IPhO2026Problems.Problem4A5.thermalPressureCoefficient_eq_inverse_referenceTemperature}
  \uses{def:physics:IPhO_2026_4_A_5:aux006, def:physics:IPhO_2026_4_A_5:aux007, def:physics:IPhO_2026_4_A_5:aux008, def:physics:IPhO_2026_4_A_5:aux009, def:physics:IPhO_2026_4_A_5:aux011, def:physics:IPhO_2026_4_A_5:aux012, lem:physics:IPhO_2026_4_A_5:aux016}
  For a positive reference state and a genuinely heated second state, the normalized secant slope of an isochoric proportionality is 1 / T₀.
\end{lemma}
\begin{proof}
  Substitute the cited governing relations in order and use the stated positivity, orientation, and branch conditions to obtain the conclusion.
\end{proof}
% --- Archon physics formalization source end ---
